\documentclass[10pt,a4paper]{amsart}
\usepackage[margin=19mm]{geometry}
\usepackage{amsmath,amssymb,mathtools}
\usepackage[scr=rsfs,cal=euler]{mathalfa}
\usepackage{xfrac}
\usepackage[unicode,hidelinks,bookmarks=false]{hyperref}
\newcommand{\ie}{i.e.\ }
\newcommand{\cf}{cf.\ }
\newcommand{\resp}{respectively\ }
\newcommand{\Subsection}{\subsection}
\providecommand{\nobreakdash}{\nobreak\hbox{-}}
\setlength{\emergencystretch}{2em}
\multlinegap0pt
\usepackage{bm}
\usepackage{bbm}
\usepackage{dsfont}
\usepackage{xspace}
\usepackage{cancel}
\usepackage{mathtools}
\usepackage{amsthm}
\makeatletter
\@ifundefined{thm}{\newtheorem{thm}{Thm}}{}
\@ifundefined{claim}{\newtheorem{claim}{Claim}}{}
\@ifundefined{lemma}{\newtheorem{lemma}[thm]{Lemma}}{}
\@ifundefined{proposition}{\newtheorem{proposition}[thm]{Proposition}}{}
\@ifundefined{theorem}{\newtheorem{theorem}[thm]{Theorem}}{}
\makeatother
\title[New examples in the study of selectively separable spaces]{New examples in the study of selectively separable spaces}
\author{Alan Dow, Hayden Pecoraro}
\date{Source: arXiv:2509.22590v2 (CC BY 4.0)}
\begin{document}
\maketitle

\noindent\emph{Source and licence.} New examples in the study of selectively separable spaces. Version arXiv:2509.22590v2 (CC BY 4.0), distributed under the Creative Commons Attribution 4.0 International licence, as recorded in the arXiv licence metadata for this exact version and shown on the abstract page https://arxiv.org/abs/2509.22590v2. Full rights notice: licenses/arxiv-2509.22590-CC-BY-4.0.txt.\par\smallskip

\noindent\textit{Editorial scope.} Counted unit: the Theorem whose source label is \texttt{None} in the article (source0.tex lines 901--904, source label \texttt{None}), delivered together with the article's own complete original proof of it (source0.tex lines 906--1183, one \texttt{proof} environment of 278 lines). No mathematics is written, repaired or re-derived: statement and proof are the article's own text line for line. Required context retained uncounted: getModels (lines 202--207); ringRule (lines 279--297); mainLemma (lines 325--350); notH (lines 474--485); justwH (lines 635--638). Every definition, display and label of those lines is retained unchanged. Omitted from this package and not redistributed: 1--201, 208--278, 298--324, 351--473, 486--634, 639--900, 905--905, 1184--1348. Nothing inside a retained excerpt is omitted or altered: every retained body is delivered exactly as the article's lines are written, so no omission receipt of any kind accompanies this package. Citations: the retained excerpts carry no citation command at all, so no bibliography entry of the article is transcribed and no reference is redistributed. Reference adaptation: none was needed and none was made; every reference inside the delivered unit resolves inside it, so no unresolved reference or citation is produced. Printed slips kept literal: no printed word, symbol, hypothesis or number of the counted statement or of its proof was altered. Layout and class: the article's own class is \texttt{amsart} with packages including xcolor, amssymb, amsrefs, hyperref; the wrapper is the lane's pinned \texttt{amsart} editorial wrapper at 10pt on A4, which restates the theorem-like environments the retained text needs and adds the article's own macro definitions that the retained text needs (none), with extra wrapper packages (bm, bbm, dsfont, xspace, cancel, mathtools). The delivered numbering is the unit's own. Assets: no retained span contains a figure environment, a graphics-inclusion command, a PDF-page inclusion, a table or a TikZ picture; no image file is copied into this package, the package declares no asset dependency and no image file is redistributed with it. Licence: version arXiv:2509.22590v2 of this article is distributed under the Creative Commons Attribution 4.0 International licence, as recorded in the arXiv licence metadata for this exact version and shown on the abstract page https://arxiv.org/abs/2509.22590v2. A full rights notice is delivered at licenses/arxiv-2509.22590-CC-BY-4.0.txt. Characters: every retained excerpt is pure ASCII, so the delivered engine, XeLaTeX with its default Latin Modern text font, reports no missing character.
\begin{proposition} For any regular cardinal $\kappa\leq \mathfrak c$
and\label{getModels}
any $S\subset H(\mathfrak c^+)$ of cardinality less than $\kappa$,
there is an elementary submodel $M$ of $H(\mathfrak c^+)$ satisfying
that $S\subset M$, $|M|<\kappa$, and $M\cap \kappa$ is an element of $\kappa$.
\end{proposition}

   \begin{proposition} 
   Assume    that $\varphi$ is a\label{ringRule}  
   function from $I(q,f)$ to $2^{<\omega}$
and let $ \tau\supset\tau_0$ be a clopen base 
for a topology on $\mathbb Q$. Assume that $q$ is in 
the closure of some $A\subset I(q,f)$ with respect to
$\tau$. Let $\sigma\supset  \tau$ be another
clopen base for a topology on $\mathbb Q$ that satisfies
    for each $q\in W\in \sigma$, there is a finite
   subset $\{ \rho_i : i\in\ell\}\subset 2^\omega$ 
   and a $U\in  \tau $,
   $$W\cap I(q,f) =^* U\cap I(q,f) \setminus
     \{ \varphi^{-1}(s) : s\in \{ \rho_i \restriction n: i\in\ell, \ n\in\omega\}~\}~.$$
Then  $q$ is in the $\sigma$-closure of $A$ unless there is some $U\in\tau$
such that 
$T_{A\cap U} = \{  s\in 2^{<\omega} :
 (\exists a \in A\cap U) s \leq \varphi(a)  \}$ has only finitely many
infinite branches.
   \end{proposition}

\begin{lemma} Suppose that $\tilde \tau \supset \tau_0$ is a
clopen base\label{mainLemma} for a topology on $\mathbb Q$ and let $\mathcal I$
be a family of $\tilde \tau $-converging sequences with $
|\tilde \tau | + |\mathcal I|<\mathfrak b$. Let $\tilde f$ be any
member of $\omega^\omega$.

Let $B\subset \mathbb Q$ be any set that is almost disjoint from each $I\in \mathcal I$.
Also let $A$ be any subset of $\mathbb Q$ and suppose that $r\in \mathbb Q\setminus A$
is not the limit of any  $\tilde \tau $-converging sequence contained in $A$.
 \medskip

 Then there is a function $\tilde f\leq f\in \omega^\omega$ and a
 clopen base $\sigma\supset \tilde \tau$ such that
 \begin{enumerate}
 \item   $I\subset^* I(q,f)$ for each $q\in \mathbb Q$ and $I\in\mathcal I$ that converges
 to $q$,
 \item $\{U_n: n\in\omega\} = \sigma\setminus \tilde \tau$ is a countable partition of $\mathbb Q$,
 \item every element of $\mathcal I$ is a $\sigma$-converging sequence,
 \item $B$ is closed and discrete with respect to $\sigma$,
 \item $r\in U_0$  and $U_0$ is disjoint from  $A$,
 \item for each $0<n\in\omega$ and $ q\in   U_n$,  
    $U_n\cap I(q,f) =^* I(q,f)\setminus B$,
  \item for $q\in U_0 $, 
   $U_0\cap I(q,f) =^* I(q,f) \setminus (B\cup A^{(1)})$.
 \end{enumerate}
 \end{lemma}

\begin{lemma}  Assume that $\tau_{\mathfrak b}\supset \tau_0$
is a clopen basis\label{notH} for a topology on $\mathbb Q$ that satisfies
\begin{enumerate}
\item  for every $t\in 2^{<\omega}$, $Q_t$ is a $\tau_{\mathfrak b}$-dense
 set,
 \item for every $\alpha\in \mathfrak b$, there is a $\rho_\alpha\in 2^\omega$
 such that $B_\alpha = \bigcup\{
  Q_{\rho\restriction n}\cap \{q_\ell : \ell < f_\alpha(n) \}\}$ 
  is closed and discrete,
\end{enumerate}
then $(\mathbb Q, \tau_{\mathfrak b})$ is not H-separable.
\end{lemma}

\begin{theorem}
There is a countable space that is mH-separable\label{justwH}
but not H-separable.
\end{theorem}

\begin{theorem}[$\mathfrak p = \mathfrak b$]
There is a countable Fr\'echet-Urysohn space that is
 not H-separable.
\end{theorem}

\begin{proof} For the reader's convenience we repeat
the main details but also introduce some minor changes at limit steps.
We recursively  construct
our  increasing family $\{ \tau_\alpha : \alpha < \mathfrak b\}$
of clopen bases for topologies on $\mathbb Q$ together with 
the special  family
 $\{ \rho_{\alpha+1} : \alpha < \mathfrak b\}
 \subset 2^\omega$ 
 by also constructing an increasing chain,
  $\{ M_\alpha : \alpha < \mathfrak b\}$,
 of elementary
 submodels of $H(\mathfrak c^+)$ as per Proposition \ref{getModels}.
 For this construction, we choose to make $\{ M_\alpha :\alpha <\mathfrak b\}$
  a continuous increasing chain,  and for limit $\alpha$,
    $\tau_\alpha $ will equal $\bigcup\{ \tau_\beta : \beta < \alpha\}$. 
Similarly for limit $\alpha$ our choice for $B_\alpha$ will be the empty set.

 For each $\alpha <\mathfrak b$,  
  let $\delta_\alpha$ be the ordinal
    $M_\alpha\cap \mathfrak b$  and again choose $\hat f_{\delta_\alpha}
    \geq f_{\delta_\alpha}$
    in $M_{\alpha+1}\cap \omega^\omega$ so as to 
    mod finite dominate $M_\alpha\cap \omega^\omega$.
  The elementarity will be
    critical for ensuring that each $I(q,\hat f_{\delta_\alpha})$ 
    captures enough information,
    combined with inductive hypotheses motivated by
    Proposition \ref{ringRule},
    to ensure that the final
    topology is Fr\'echet-Urysohn.

For $0<\gamma  < \alpha$
 and $\beta+1 < \alpha$ we assume we are preserving
these inductive properties:
   
\begin{enumerate}
\item for each $t\in 2^{<\omega}$, $Q_t$ is $\tau_{\gamma}$-dense,
\item $M_\gamma$ is an elementary submodel  of $H(\mathfrak c^+)$
satisfying that $|M_\gamma|<\mathfrak b$ and $M_\gamma\cap \mathfrak b = 
 \delta_\gamma\in \mathfrak b$,
 \item each of $ \{Q_t:t\in 2^{<\omega}\}$ and $\{f_\xi: \xi\in\mathfrak b\}$
  are elements of $M_0$,
\item  $\{ M_\xi : \xi\leq \beta\}$ is an  element  of $M_{\beta+1}$,
\item $\hat f_{\delta_\beta}$ is an element of $M_{\beta+1}\cap \omega^\omega$
that is strictly increasing and mod finite dominates $M_\beta\cap \omega^\omega$,
\item   $\tau_\beta$ is a subset of $M_\beta$  and $\tau_{\beta+1}\setminus \tau_\beta\in M_{\beta+1}$ is countable,
\item $\rho_{\beta+1}\in 2^\omega\cap M_{\beta+1} \setminus M_\beta $,
 \item the\label{H3} set $B_{\beta+1} = \bigcup\{ Q_t \cap \{ q_\ell : \ell < f_{\beta}(|t|) \}: 
  t\in \{ \rho_{\beta+1}\restriction n : n\in \omega\}\}$
is closed and discrete with respect to 
$\tau_{\beta+1}$,
 \item for\label{mH1} each $q\in W\in \tau_\gamma$ and $\mu<\gamma$, there is a
 $q\in U\in \tau_\mu$
 such that $ U\cap I(q,\hat f_{\delta_\mu}) \setminus W$ is 
 covered by a finite union from $\{ B_{\xi+1} : \mu \leq \xi < \gamma\}$.
\end{enumerate}

We choose $M_0$ exactly as in the proof of
 Theorem \ref{justwH}. For a limit ordinal $\alpha < \mathfrak b$,
   we set $ M_\alpha = \bigcup\{ M_\beta : \beta < \alpha\}$ and
      $\tau_\alpha = \bigcup\{ \tau_\beta : \beta < \alpha\}$ and there
      is nothing more to do.

      Now assume that $\alpha = \beta+1$.  
We choose   any elementary submodel $M_\alpha$ of $H(\mathfrak c^+)$
such that $ \{M_\beta , \tau_\beta\}\in M_\alpha$
 and such that $|M_\alpha | < \mathfrak b$
and $M_\alpha\cap \mathfrak b =\delta_\alpha>\alpha$. 
We also choose $\hat f_{\delta_\beta}\geq f_{\delta_\beta}$
to be any element
of $M_\alpha\cap \omega^\omega$ that is strictly increasing
and that mod finite dominates
 $M_\beta\cap \omega^\omega$.
Choose any $\rho_\alpha \in  2^\omega\cap M_\alpha \setminus M_\beta$ 
and set  $B_\alpha =
\bigcup\{ Q_t \cap \{ q_\ell : \ell < f_\alpha(|t|) \}: 
  t\in \{ \rho_\alpha\restriction n : n\in \omega\}\}$.
  The set $B_\alpha$ is an element of $M_\alpha$.  
  Define
  $\tau_\alpha$ to be $\tau_\beta \cup
  \{U^\alpha_n : n\in \omega\}$ where $\{U^\alpha_n : n\in\omega\}$
  is obtained by applying Lemma \ref{mainLemma}, working within
   $M_\alpha$, to the
  values $\tilde\tau = \tau_\beta$,
  $\mathcal I = \emptyset$,
   $A=\emptyset$, $r=0$, $\tilde f = \hat f_{\delta_\beta}+f_\alpha  $, and
    $B = B_\alpha$.  
    
     
This completes the construction, and we set $\tau_{\mathfrak b}$
equal to $\bigcup\{ \tau_\alpha : \alpha < \mathfrak b\}$.  
 It follows by Lemma \ref{notH} that $\tau_{\mathfrak b}$ is not
  H-separable. Now
we prove that $(\mathbb Q, \tau_{\mathfrak b})$ is 
 Fr\'echet-Urysohn.  
 
 
 Consider any $A\subset \mathbb Q$
 and $\tau_{\mathfrak b}$-limit point $q$ of $A$.   
 We again note that $q$ is a $\tau_\alpha$-limit point 
 of $A$ for every $\alpha < \mathfrak b$. 
Fix a continuous increasing chain, $\{ M^q_\alpha : \alpha < \mathfrak b\}$,
   of elementary submodels of $H(\mathfrak c^+)$ satisfying that
   $A$ and 
     $\{ M_\beta: \beta < \mathfrak b\}$ are elements of $ M^q_0$ and for all
      $\alpha<\mathfrak b$, $M^q_\alpha$ has cardinality less
      than $\mathfrak b$ and $\mu_\alpha = M^q_\alpha\cap \mathfrak b$
      is an element of $\mathfrak b$.  Let $C$ be a closed 
      and unbounded set of $\zeta\in \mathfrak b$ that satisfy
       $M^q_\zeta \cap \mathfrak b = \zeta$. Choose a strictly
       increasing sequence $\{ \zeta_k : k\in\omega\}\subset C$,
        and let $\mu$ be the supremum. Note also that, for $k\in\omega$,
         $M^q_{\zeta_k}\subset M^q_\mu$ and that 
        $\delta_\mu = \mu$.

For each $k\in\omega$, using that $\tau_{\zeta_k}$ is a
 Fr\'echet-Urysohn topology, choose a sequence $J_k\subset A$
 that is an element of $M^q_{\zeta_{k+1}}$ and that 
 $\tau_{\zeta_k}$-converges to $q$.    
Let  $L_k = \{ n\in  \omega : J_k\cap U(q,n)\setminus U(q,n+1)\neq\emptyset\}$
 and define the  function 
 $g_k\in\omega^\omega\cap M^q_{\zeta_{k+1}}$ 
 according to the rule that for each
  $n\in L_k$, 
  there is an $\ell< g_k(n)$ such
that $q_\ell \in I_k\cap U(q,n)\setminus U(q,n+1)$
and for each  $m\in\omega\setminus L_k$,
  $g_k(m) = g_k(\min(L_k\setminus m))$.  
By elementarity, there is a $\xi_k\in \mathfrak b
\cap M^q_{\zeta_{k+1}}$ such that $f_{\xi_k}(n)> g_k(n)$
for infinitely many $n$. It follows that $\hat f_{\delta_{\zeta_{k+1}}}(n)
> g_k(n)$ for infinitely many $n$. Since $\hat f_{\delta_{\zeta_{k+1}}}$ is
a strictly increasing function, it follows that for each $n$
such that $\hat f_{\delta_{\zeta_{k+1}}}(n)
> g_k(n)$, we also have that $g_k(\min(L_k\setminus n)) < 
\hat f_{\delta_{\zeta_{k+1}}}(\min(L_k\setminus n))$. 
This implies that 
$I(q,\hat f_{\delta_{\zeta_{k+1}}})\cap J_k$ is infinite.
Since $\delta_{\zeta_{k+1}}<\delta_\mu=\mu$,
 it similarly follows that, for each $k\in\omega$,
    $I(q,\hat f_{\mu}) \cap J_k$ is infinite.  

    \begin{claim}  $q$ is in the $\tau_\mu$-closure of
       $A\cap I(q,\hat f_{\mu})$. 
    \end{claim}

  \bgroup


\def\proofname{Proof of Claim.}

\begin{proof} Let $U$ be any element of $\tau_\mu$. 
  Choose $k\in\omega$ so that $U\in \tau_{\zeta_k}$
   and observe that  since $U$ contains a cofinite subset of $J_k$,
  $   U\cap (J_k\cap I(q,\hat f_{\mu}))\subset U\cap A\cap I(q,\hat f_\mu)$ is infinite.
\end{proof}


Since $|\tau_\mu|<\mathfrak p$ there are sequences contained
 in $A\cap I(q,\hat f_\mu)$ that $\tau_\mu$-converge to $q$. Now fix any $\beta\in \mu$
 and, trivially, observe that ``there exists some $\mu>\beta$ satisfying
 that there exists some sequence contained in $A\cap I(q,f_{\delta_\mu})$
 that  $\tau_{\mu}$-converges to $q$''. 
 Now, working in $M^q_\mu$, apply elementarity,
  and choose, for each $k\in\omega$, some $\zeta_k < \mu_k < \mu$ such
  that there is sequence $I_k\subset A\cap I(q,\hat f_{\delta_{\mu_k}})$ 
  that $\tau_{\mu_k}$-converges to $q$. By passing to subsequences,
   we can assume that, for each $k$, 
     $\varphi(I_k)$ is one of: a singleton $t_k\in 2^{<\omega}$,
      an infinite anti-chain $T_k\subset 2^{<\omega}$, or
      an infinite subset of $\{ \psi_k\restriction n : n\in\omega\}$
      for some $\psi_k\in 2^\omega\cap M^q_\mu$. 

 \begin{claim} If, for some $k\in\omega$, $\varphi(I_k)$ is either a singleton
  or an anti-chain, then $I_k$ converges to $q$ with respect to $\tau_{\mathfrak b}$. 
 \end{claim}


 \begin{proof}  
 Consider any $W\in \tau_{\mathfrak b}$.  By inductive condition \ref{mH1},
  there is   $U\in \tau_{\mu_k}$ and a finite set $\{\xi_i : i<\ell\}\subset
     \mathfrak b \setminus \mu_k$ satisfying that
       $U\cap I(q,\hat f_{\mu_k}) \setminus W$ is contained in
         $\bigcup\{ B_{\xi_i+1} : i< \ell\}$.  
Of course $\bigcup\{ \varphi(B_{\xi_i+1}) : i<\ell\}$ is a subset of
   $\bigcup \{ \rho_{\xi_i+1}\restriction n : i< \ell,\ \ n\in\omega\}$. 
The set $\bigcup\{ B_{\xi_i+1} : i< \ell\}$ is almost disjoint
from every $Q_t$ ($t\in 2^{<\omega}$) and for each
 $t\in 2^{<\omega}\setminus 
  \bigcup \{ \rho_{\xi_i+1}\restriction n : i< \ell,\ \ n\in\omega\}$,
  it 
      is disjoint from
      $Q_t$.   
      It therefore follows that $I_k\cap 
        \bigcup\{ B_{\xi_i+1} : i< \ell\} $ is finite. 
        Since $U$ contains $I_k$ mod finite, it follows
        that $W$ does as well.
 \end{proof}

 The remaining case is that, for each $k$, $\varphi(I_k)$
 is an infinite chain contained in $\{ \psi_k\restriction n : n\in \omega\}$
 with $\psi_k\in M^q_\mu$.  
 
 \begin{claim} If for some $k\in\omega$, $q$  is in the $\tau_{\mu}$-closure
 of $I_k$, then $I_k$ has a subsequence that $\tau_{\mathfrak b}$-converges to $q$.
 \end{claim}

 \begin{proof}  If $\psi_k\in \{ \rho_{\xi+1} : \xi <\mathfrak b\}$, then,
  by elementarity, 
 choose $\xi_k\in M^q_\mu$ so that $\psi_k = \rho_{\xi_k+1}$.  Choose 
  $ k +1 < \bar k\in \omega$ so that $\xi_k < \mu_{\bar k}$. 
  By assumption $I_k \subset I(q, \hat f_{\delta_{\mu_{\bar k}}})$ and
     $q$ is a $\tau_{\mu_{\bar k}}$-limit of $I_k$. Choose any infinite
      $I'_k\subset I_k$ that converges to $q$ with respect to 
      $\tau_{\mu_{\bar k}}$. Note that $I'_k$ is almost disjoint
      from $B_{\gamma+1}$ for all $\mu_{\bar k}\leq \gamma$. 
      Therefore, by induction hypothesis \ref{mH1}, $I_k'$ 
      converges to $q$ with respect to $\tau_{\mathfrak b}$. 
 \end{proof}


  Therefore, we now assume that $q$ is not in the $\tau_\mu$-closure
  of $I_k$ for each $k\in \omega$, and, as a consequence we obtain
  this next claim.

  \begin{claim} The
  set $\{ \psi_k : k\in\omega\}$ is infinite.
  \end{claim}

  \begin{proof} Assume that there is a single $\psi\in 2^\omega$
  such that $\psi_k = \psi$ for infinitely many $k\in \omega$. 
  As in the previous proof, choose $\bar k$ so that, if 
   $\psi \in \{\rho_\xi : \xi\in\mathfrak b\}$, then
      $\psi \in \{ \rho_\xi : \xi < \mu_{\bar k}\}$. 
   Choose any $k >\bar k$ so that $\psi = \psi_k$. Now, 
   we have that for all $\mu_{ k} < \beta < \mu$,
   $U^\beta_{m_q} $ mod finite contains
    $ I(q,\hat f_{\delta_\beta}) \setminus B_{\beta+1}$, 
    where $m_q$ is chosen so that $q\in U^\beta_{m_q}$. 
    Since $I_k$ is mod finite contained in 
       $I(q,\hat f_{\delta_\beta})\setminus B_{\beta+1}$
       (because $\rho_{\beta+1}\neq \psi$) 
       it would follow that $I_k$ converges to $q$ with 
        respect to $\tau_\mu$.
  \end{proof}

By the previous claim we can, by re-indexing an infinite subsequence
 of $\{ \psi_k : k\in \omega\}$, assume that they are pairwise
 distinct and converge to some $\psi\in 2^\omega$ and that
    $\{ n_k : k\in\omega\}$ strictly increases where
     $n_k = \min\{ n : \psi(n)\neq \psi_k(n)\}$. 
  Next, we pass
 to a cofinite subset of each $I_k$ and thereby assume that
 the sets  $\{ \varphi(I_k) : k\in \omega\}$ are pairwise disjoint
 by simply ensuring that $|t|>n_k$ for all $ t\in \varphi(I_k)$.


 \begin{claim} For every $U\in \tau_\mu$, $I_k\subset^* U$ for all
 but finitely many $k\in \omega$.
 \end{claim}

 \begin{proof}  For each $U\in \tau_{\mu_k}$, $I_\ell\subset ^* U$
 for all $k<\ell$ since $I_\ell$ converges to $q$ with respect
 to $\tau_{\mu_\ell}$ and $U\in \tau_{\mu_\ell}$. 
 \end{proof}


  \egroup

  Finally, since $|\tau_\mu|<\mathfrak b$, there is a selection 
   $S\subset \bigcup\{ I_k : k\in \omega\}$  ($|S\cap I_k| = 1$
   for all $k\in \omega$)  satisfying that 
    $S\subset^* U$ for all $U\in \tau_\mu$.
We have produced a sequence $S$ that converges to $q$ with respect
to $\tau_\mu$ such that $\varphi(S)$ is an infinite anti-chain.
Clearly $B_{\beta+1}\cap S$ is finite for all $\mu\leq\beta<\mathfrak b$,
 and, by induction hypothesis \ref{mH1}, this completes the proof.
\end{proof}

\end{document}
